Sauerstoff (O$_2$) ist zweiatomig. Bei dieser Temperatur sind die drei Translations- und die zwei Rotationsfreiheitsgrade angeregt, also $f = \fOO$. Die Stoffmenge ist $n = m/M$ mit $M = \MOO$ (FoTaBe S.~228):
\begin{align}
 U &= \frac{f}{2}\,nRT = \UOxF \\
   &= \frac{\fOO}{2}\cdot\frac{\mO}{\MO}\cdot\ncR\cdot\TO \\
   &= \UOx \approx \result{\UOxP{3}{2}}
\end{align}
